
\chapter{Kurzfassung}

\emph{Beat the Score} ist ein Trainingsprogramm für beginnende Musiker. Es hilft dem Benutzer beim Üben die richtigen Noten zum richtigen Zeitpunkt zu treffen. Es erfolgt eine sofortige Rückmeldung darüber wie der Benutzer gespielt hat. Der Computer begleitet den Spieler dabei mit weiteren Instrumenten.

Durch Unterstützung des populären GP5 Formats kann der Benutzer eine große Auswahl an Noten im Internet finden.
Für Rückmeldung der Spielweise können Instrumente per Klinke oder USB an den Computer angeschlossen werden. Unterstützt werden Gitarren und Klaviere.

Da Anfänger oft nicht Notenschrift lesen können, werden verschiedene Anzeigemethoden angeboten: 
\begin{description}
	\item[Klaviatur] \hfil \\ Die Noten werden als Rechtecke angezeigt, wobei eine Achse die Tonhöhe und die andere den Zeitpunkt darstellt.
	\item[Notenschrift] \hfil \\ Wie sie unter Musikern üblich ist.
	\item[Tabulatur] \hfil \\ Wie sie von Gitarristen verwendet wird.
\end{description}

Zum Üben ist es auch möglich, das Tempo des Liedes einzustellen oder den Trainingsmodus zu aktivieren, wo bei jeder Note pausiert wird, bis die richtige Note gespielt wurde.

Das Programm ist erhältlich im Google Play Store für Android und auf der Website {\tt http://beatthescore.at} für Linux, Windows und in Zukunft auch Mac OSX.